\section{Whitney Topologies}

Let $X,Y$ be smooth manifolds. In this section we discuss the concept of $C^k$-convergence for sequences of functions $(f_n) \subseteq C^k(X,Y)$ with respect to the Whitney topolgies. The notation and terminology used will be the same as in \cite{golubitsky2012stable}. First we remark that $X$ or $Y$ is also allowed to have boundary. The definition of the jet bundles $J^k(X,Y)$ and the Whitney topologies works exactly the same (if both have boundary, $J^k(X,Y)$ will be a manifold with corners, so we might want to avoid that). Let us begin by understanding the $C^0$-topology.

\begin{lemma}
    Let $X,Y$ be smooth manifolds with $X$ compact. Then the Whitney $C^0$-topology (as described in \cite{golubitsky2012stable}) and the exponential topology (as described in \cite{steimle_topo}) on $C^0(X,Y)$ agree.
\end{lemma}

\begin{remark}
    If $X$ is not compact the $C^0$-topology will in general be strictly finer than the exponential topology. To see this recall that a subbase for the exponential topology on $C^0(X,Y)$ is given by sets of the form
    \[
    F(K,V) := \{f \in C^0(X,Y) \mid f(K) \subseteq V\}
    \]
    for compact sets $K \subseteq X$ and open sets $V \subseteq Y$. Set 
    \[
    W := X \times V \cup (X\setminus K) \times Y.
    \]
    Then $W \subseteq X \times Y$ is open and 
    \[
    F(K,V) = M(W).
    \]
    To see that the $C^0$ is finer in general consider the following example. ??????????????
    To differentiate between the two topologies on $C^0(X,Y)$ we shall use the term $C^0_\mathrm{exp}$-convergence to refer to convergence with respect to the exponential topology.
\end{remark}

Next we understand convergence in the $C^k$-topology. Before stating the main proposition we shall introduce another type of convergence which is better suited for non-compact domains, since in the non-compact case usual $C^k$-convergence is often too strong, because 
\[
f_n \xrightarrow{C^k} f
\]
necessarily requires for $f_n$ and $f$ to agree outside a compact subset (see \cite{golubitsky2012stable}).

\begin{definition}
    Let $X,Y$ be smooth manifolds and $f,f_n \in C^k(X,Y)$. We say that 
    \[
    f_n \xrightarrow{C^k_{\mathrm{loc}}} f
    \]
    if for every compact codimension $0$-submanifold $K \subseteq X$ we have 
    \[
    f_n|_K \xrightarrow{C^k} f|_K.
    \]
\end{definition}

Note that for compact domains $C^k_{\mathrm{loc}}$-convergence and $C^k$-convergence coincide, so we are actually interested in understanding $C^k_{\mathrm{loc}}$-convergence. The following two propositions give a useful characterization of $C^k_{\mathrm{loc}}$-convergence to work with in practice. 

\begin{proposition}
    \label{whitney_top:prop:ck_conv_on_rn}
    Let $D\subseteq \R^n$ be a compact codimension $0$ submanifold, $V \subseteq \R^m$ open and $f,f_n \in C^k(D,V)$. Then
    \[
    f_n \xrightarrow{C^k} f
    \]
    if and only if for every multiindex $\alpha \in \N^n_0$ with $0\leq\abs{\alpha} \leq k$
    \[
    \frac{\partial^{\abs{\alpha}}f_n}{\partial x^\alpha} \xrightarrow{\mathrm{uniformly}}\frac{\partial^{\abs{\alpha}}f}{\partial x^\alpha}.
    \]
\end{proposition}

\begin{proposition}
    Let $X,Y$ be smooth manifolds and $f,f_n \in C^k(X,Y)$. Then the following statements are equivalent.
    \begin{enumerate}[(i)]
        \item $f_n \xrightarrow{C^k_{\mathrm{loc}}} f$,
        \item $j^kf_n \xrightarrow{C^0_{\mathrm{loc}}} j^kf$,
        \item For all charts $\phi \colon U \to U^\prime \subseteq X$ and $\psi \colon V \to V^\prime \subseteq Y$ with $f(U^\prime), f_n(U^\prime) \subseteq V^\prime$ (for $n$ sufficiently large) we have 
        \[
        \psi^{-1} \circ f_n \circ \phi \xrightarrow{C^k_{\mathrm{loc}}} \psi^{-1} \circ f \circ \phi,
        \]
        \item For every $p \in X$ there exist charts $\phi \colon U \to U^\prime \subseteq X$ around $p$ and $\psi \colon V \to V^\prime \subseteq Y$ around $f(p)$ with $f(U^\prime), f_n(U^\prime) \subseteq V^\prime$ (for $n$ sufficiently large) such that
        \[
        \psi^{-1} \circ f_n \circ \phi \xrightarrow{C^k_{\mathrm{loc}}} \psi^{-1} \circ f \circ \phi.
        \]
    \end{enumerate}
\end{proposition}

Before proving the propositions we recall a few facts about the exponential topology not mentioned in \cite{steimle_topo}, which will be used in the proofs.

\begin{lemma}
    Let $X,Y,Z$ be topological spaces and $A \subseteq X$, $B \subseteq Y$ subspaces. Then the following holds (assuming the exponential topology exists on each space)
    \begin{enumerate}[(i)]
        \item The canonical map $C^0(X,B) \hookrightarrow C^0(X,Y)$ is an embedding
        \item The restriction map $C^0(X,Y) \to C^0(A,Y)$ is continuous
        \item The canonical bijection $C^0(X,Y) \times C^0(X,Z) \to C^0(X,Y \times Z)$ is a homeomorphism
    \end{enumerate}
\end{lemma}

\begin{proof}
    For (i) we show that the exponential topology on $C^0(X,B)$ agrees with the topology induced by the map  $C^0(X,B) \hookrightarrow C^0(X,Y)$ by showing the corresponding universal property. So let
    \[\begin{tikzcd}
	& {C^0(X,Y)} \\
	Z & {C^0(X,B)}
	\arrow["{F}", from=2-1, to=1-2]
	\arrow["f", from=2-1, to=2-2]
	\arrow[hook, from=2-2, to=1-2]
\end{tikzcd}\]
commute with $F$ being continuous. We have to show that $f$ is continuous, but by definiton of the exponential topology this is equivalent to 
\[
Z \times X \to B, \; (z,x) \mapsto f(z)(x)
\]
being continuous and continuity of $F$ is equivalent to continuity of 
\[
Z \times X \to Y, \; (z,x) \mapsto F(z)(x),
\]
so (i) follows. For (ii) note that the restriction map is continuous if and only if 
\[
C^0(X,Y) \times A \to Y, \; (f,x) \mapsto f|_A(x)
\]
is continuous. But this is just the restriction of the evaluation map which is continuous by \cite{steimle_topo}. Similiarly the map in (iii) is continuous if and only if 
\[
C^0(X,Y) \times C^0(X,Z) \times X \to Y \times Z, \; (f,g,x) \mapsto (f(x),g(x))
\]
is. But again the continuity follows from the continuity of the evaluation map.
\end{proof}

\begin{proof}[Proof of Proposition \ref{whitney_top:prop:ck_conv_on_rn}]
    In \cite{golubitsky2012stable} it was shown that for compact domains $C^k$-convergence is equivalent to the $k$-jets converging uniformly with respect to any metric on the Jet-bundle compatible with its topology, that is 
    \[
    f_n \xrightarrow{C^k} f \; \Leftrightarrow \; j^kf_n \xrightarrow{\mathrm{uniformly}} j^kf
    \]
    Note that for compact domains uniform convergence and $C^0_\mathrm{exp}$-convergence agree (see \cite{steimle_topo}). Now let $d(n,k) = \frac{(n+k)!}{n!k!}$ be the total number of all partial derivatives $\partial^\alpha$ of order $1 \leq \abs{\alpha} \leq k$ and let 
    \[
    \tau \colon D \times V \times \R^{d(n,k)} \to J^k(D,V)
    \]
    be the canonical diffeomorphism. Then 
    \[
    j^kf_n \xrightarrow{C^0_\mathrm{exp}} j^kf \; \Leftrightarrow \; j^kf_n \circ \tau \xrightarrow{C^0_\mathrm{exp}} j^kf \circ \tau.
    \]
    But a sequence of maps with values in a product $C^0_\mathrm{exp}$-converges if and only if each component $C^0_\mathrm{exp}$-converges. Since the components of $j^kf_n \circ \tau$ and $j^kf \circ \tau$ are precisely the partial derivatives of $f_n$ and $f$ up to order $k$ the statement follows (after again applying the equivalence of uniform and $C^0_\mathrm{exp}$-convergence on compact domains).
\end{proof}